\honest{Science Isn't Strict Enough}{Eliezer Yudkowsky}{2008}

Once, at eighteen, I had a stupid theory. I made sure it had experimental consequences and said I wished to test it, which was all the traditional ideal of Science asked of me. It was nowhere near the effort needed to get it right.

The ideals of Science, I say, ``too readily give out gold stars.'' Test your theory, accept the result, and you are a good scientist, right or not. I mean the ideal ``as it is traditionally preached,'' not the practice of science, and I quote no one who preaches it. \nb{A commenter, poke, replied that this ideal is ``a second-hand imprecise (idealized) description of science.''}

Science began as a rebellion against armchair reasoning, so it has no rule about which hypotheses to test; that is left to the individual. It asks only for ``the ability to accept reality when you're beat over the head with it,'' and that is enough to sustain a scientific culture. Science, I say, ``wouldn't try to give an official verdict on the best hypothesis to test, in advance of the experiment.'' \nb{Karl Popper, the tradition's best-known modern philosopher, did give one: prefer the theory with more content, which is the less probable theory. So both standards rank hypotheses before the test; they rank by different things.}

Probability theory is stricter. If 1\% of screened women have breast cancer, 80\% of those with cancer test positive and 10\% of those without do, then a woman who tests positive has cancer with probability 7.5\%, ``not 7.4\% or 7.6\%.'' I quote the third virtue, lightness, from my own ``Twelve Virtues of Rationality.''

Kahneman, Tversky and Jaynes taught me a stricter way, one that could have told me in advance: ``You picked the wrong hypothesis to test, dunderhead.'' I grant that it is much harder to use, that one error in a hundred steps can carry you anywhere, and that Science has good reason not to trust you with it. But if you do not want to waste ten years, you must try.

Not knowing the priors, I say, ``doesn't mean you get a free, personal choice of making the priors whatever you want.'' If whim could make priors correct, you could know things without looking. \nb{Here Yudkowsky takes one side of a dispute among Bayesians without mentioning it. Subjective Bayesians such as de Finetti and Savage hold that every coherent prior is permitted. They do not claim that such a prior is accurate, which is the claim the argument answers.} Bayes, I add, is ``the law'' behind every statistical tool.

I end with Nick Tarleton's comment, which says it better than I did: ``The problem is encouraging a private, epistemic standard as lax as the social one.''
